\chapter{Bleeding Between Periods (Spotting)}

\section*{Definition}
Vaginal bleeding or spotting outside the expected menstrual period. Bleeding after menopause is a separate presentation that always needs medical assessment.

\section*{When to See a Doctor}

\begin{itemize}
 \item Arrange assessment for any unexplained bleeding between periods, especially if it is recurrent, heavy, associated with pain or discharge, or occurs after menopause.
 \item Seek urgent care for heavy bleeding, faintness, severe pelvic pain, or bleeding after a missed period, because ectopic pregnancy must be excluded.
\end{itemize}

\section*{How It Develops}
Intermenstrual bleeding may arise from hormonal contraception or ovulatory changes, infection, pregnancy, cervical disease, or uterine conditions such as polyps, fibroids, hyperplasia, or cancer. Bleeding after menopause can have benign causes but always requires evaluation to exclude endometrial or other genital-tract disease.

\section*{Causes}
\begin{itemize}
 \item Hormonal imbalance (anovulatory cycles, perimenopause)
 \item Uterine fibroids or polyps
 \item Cervical polyps, cervicitis, or ectropion
 \item Endometrial hyperplasia or carcinoma
 \item Hormonal contraception (IUD, pills, implant) breakthrough
 \item Pregnancy complications (ectopic pregnancy, threatened miscarriage)
 \item Infection (pelvic inflammatory disease, sexually transmitted infections)
 \item Thyroid disorders
 \item Coagulation disorders
 \item Cervical or vaginal trauma
\end{itemize}

\section*{Related Symptoms}
\begin{itemize}
 \item Abnormal menstrual cycles
 \item Pelvic pain
 \item Fatigue
 \item Heavy periods
 \item Vaginal discharge
\end{itemize}

\section*{Medical Evaluation}
\begin{itemize}
 \item Your doctor will take a detailed menstrual and reproductive history and perform a pelvic examination including speculum and bimanual evaluation.
 \item Transvaginal ultrasound assesses endometrial thickness, ovarian morphology, and identifies fibroids, polyps, or adnexal masses.
 \item Testing usually includes a pregnancy test when pregnancy is possible and a blood count when bleeding is heavy or prolonged. Thyroid or other tests are guided by symptoms and risk factors; prolactin is not routine for all bleeding.
 \item Endometrial sampling is considered for persistent abnormal bleeding, age 45 or older, or younger patients with risk factors such as prolonged anovulation, obesity, or failed treatment. The method and need for hysteroscopy depend on local guidance and findings.
 \item Ultrasound, hysteroscopy, or specialist referral is selected according to the history, examination, age, pregnancy status, and suspected cause.
\end{itemize}

\section*{Treatment Options}
\begin{itemize}
 \item Management depends on the cause. Contraceptive-related bleeding may improve with time or a clinician-guided change in method; hormonal treatment is selected after pregnancy and other causes are considered.
 \item Endometrial polyps and submucosal fibroids are managed with hysteroscopic resection.
 \item Endometrial hyperplasia without atypia is treated with progestin therapy (oral or intrauterine) and followed with surveillance biopsies.
 \item Atypical hyperplasia and endometrial cancer require specialist management; surgery may be recommended based on pathology, stage, fertility wishes, and overall health. Hysterectomy is not an automatic treatment for unexplained bleeding.
\end{itemize}

\section*{Self-Care and Home Management}
\begin{itemize}
 \item Track the bleeding pattern in a menstrual diary: note the date, duration, amount (spotting vs. heavy), and any associated symptoms like pain or cramping.
 \item Use the menstrual products you normally tolerate safely; change tampons as directed and never leave them in longer than the product instructions allow.
 \item If bleeding is heavy and a clinician has confirmed that NSAIDs are safe for you, they may reduce menstrual blood loss and cramps, but they do not treat unexplained intermenstrual bleeding. Avoid them when contraindicated.
 \item Take a home pregnancy test if there is any possibility of pregnancy, as implantation bleeding and ectopic pregnancy can present as spotting.
\end{itemize}

\section*{References}
\begin{thebibliography}{99}
\bibitem{bleeding_between_periods:ref1} Munro MG, Critchley HOD, et al. ``FIGO classification of causes of abnormal uterine bleeding.'' \textit{Int J Gynaecol Obstet}. 2020;149(1):3-18.
\bibitem{bleeding_between_periods:ref2} American College of Obstetricians and Gynecologists. ``ACOG practice bulletin: management of abnormal uterine bleeding.'' \textit{Obstet Gynecol}. 2023;141(5):e110-e130.
\bibitem{bleeding_between_periods:ref3} Whitaker L, Critchley HOD. ``Abnormal uterine bleeding.'' \textit{BMJ}. 2016;355:i5324.
\bibitem{bleeding_between_periods:ref4} Berek JS, Berek DL, et al. \textit{Berek \& Novak\'s Gynecology}, 16th ed. Wolters Kluwer, 2019.
\bibitem{bleeding_between_periods:ref5} American College of Obstetricians and Gynecologists. ``Abnormal Uterine Bleeding.'' ACOG patient guidance, accessed September 2026. https://www.acog.org/womens-health/faqs/abnormal-uterine-bleeding.
\bibitem{bleeding_between_periods:ref6} National Institute for Health and Care Excellence. ``Heavy menstrual bleeding: assessment and management (NG88).'' Updated 2021. https://www.nice.org.uk/guidance/ng88.
\end{thebibliography}
